% Two-page summary of the detection study in post-training/detection/.
% Build from post-training/detection/paper/: latexmk -pdf main.tex
% Figures are read from ../artifacts/figures/ (scripts/figures.py output).
\documentclass[10pt,twocolumn]{article}
\usepackage[letterpaper,margin=0.58in,columnsep=0.2in]{geometry}
\usepackage[T1]{fontenc}
\usepackage{mathptmx}
\usepackage{microtype}
\usepackage{graphicx}
\usepackage{booktabs}
\usepackage{amsmath}
\usepackage{amssymb}
\usepackage[numbers,sort&compress]{natbib}
\usepackage[hidelinks]{hyperref}
\usepackage{xcolor}
\usepackage{enumitem}
\usepackage{caption}
\usepackage{tikz}
\usetikzlibrary{positioning,fit,calc,arrows.meta,decorations.pathreplacing}
\captionsetup{font=small,labelfont=bf,skip=3pt}
\setlist{nosep,leftmargin=*}
\setlength{\bibsep}{0pt plus 0.3pt}
\graphicspath{{../artifacts/figures/}}
\renewcommand{\bibfont}{\scriptsize}
\newcommand{\ok}{\textcolor{green!45!black}{+}}
\newcommand{\no}{\textcolor{red!70!black}{$-$}}
\newcommand{\eq}{$\circ$}

\title{\vspace{-2.2em}\Large\bfseries One-Pass Object Detection from a Frozen Small
Vision-Language Model with Slots}
\author{Jakob Frick}
\date{\vspace{-2.4em}}

\begin{document}
\maketitle

\begin{abstract}
\noindent Small vision-language models (VLMs) such as Moondream detect objects by
decoding coordinates one token at a time, with one prompt per category, so latency
grows with both the number of categories and the number of objects. We post-train a
\emph{frozen} 1.9B-parameter Moondream~2 with a slot head: 48 slots of six blank
positions appended after the image tokens, filled in all at once in a single forward pass
and read through the model's own coordinate, size and vocabulary heads, trained with a
Hungarian set loss. Only low-rank adapters at slot positions and an objectness layer are
trained. On 500 COCO val2017 images the head reaches F1 0.707 at IoU~0.5 in 0.09\,s per
image (Moondream~3.1, 9B: 0.700 in 0.21\,s), and 0.718 when it also reads a $2\times2$
grid of tiles, which brings its recall of objects smaller than one feature-grid cell to
Moondream's level. On Pascal VOC 2007, which none of the models was tuned on, it scores
F1 0.858 against Moondream~3.1's 0.857. With all 80 categories requested it needs about
59$\times$ less time per image. An auxiliary masked-diffusion denoising loss, category
prompts, size-weighted losses and teacher-only labels do not help; we report what did
and did not work.
\end{abstract}

\section{Introduction}
Generalist VLMs now localise objects through text: Pix2Seq-style coordinate
tokens~\citep{chen2022pix2seq}, as in Kosmos-2~\citep{peng2023kosmos2},
Florence-2~\citep{xiao2024florence2}, PaliGemma~\citep{beyer2024paligemma} and
Moondream~\citep{moondream,moondream31}. Decoding is sequential: Moondream~2's
\texttt{detect} spends three decoder steps per box and one call per phrase. Set
prediction~\citep{carion2020detr} removes the sequence but usually needs a dedicated
detector~\citep{liu2024groundingdino,minderer2022owlvit}. We ask whether
post-training alone can give a small VLM a set output, \emph{n slots, each ``nothing'' or a
box with a label}, so that one forward pass answers every category at once. The aim is
reliable rather than exhaustive detection.


\begin{figure*}[t]
\centering
\resizebox{\textwidth}{!}{%
\begin{tikzpicture}[
  font=\footnotesize,
  base/.style={draw=black!60,fill=white,rounded corners=2pt,align=center,inner sep=3pt},
  added/.style={draw=red!70!black,fill=red!12,rounded corners=2pt,align=center,inner sep=3pt},
  outbox/.style={draw=blue!60!black,fill=blue!6,rounded corners=2pt,align=left,inner sep=5pt},
  train/.style={draw=green!40!black,fill=green!6,rounded corners=2pt,align=center,inner sep=4pt},
  arr/.style={-{Latex[length=1.8mm]},black!70,thick},
  x=1cm,y=1cm,
]
% ---------------- input
\node[font=\bfseries] at (1.0,5.85) {Input};
\draw[black!60,fill=black!4] (0.1,2.0) rectangle (1.9,4.0);
\draw[orange!80!black,thick] (0.3,2.3) rectangle (0.9,3.4);
\draw[orange!80!black,thick] (1.1,3.3) rectangle (1.6,3.8);
\node at (1.0,1.7) {image};
% ---------------- base model frame
\draw[black!50,fill=black!7,rounded corners=4pt] (2.5,0.2) rectangle (12.7,5.5);
\node[anchor=north west,font=\bfseries] at (2.6,5.45) {Base model: Moondream 2 (1.9B parameters), frozen};
\node[base] (enc) at (3.7,3.0) {vision\\encoder};
\draw[arr] (1.9,3.0) -- (enc.west);
% input sequence of the decoder: image tokens, then the added slots
\node[font=\scriptsize] at (6.0,4.62) {729 image tokens};
\foreach \i in {0,1,2} \draw[black!60,fill=white] (5.2+\i*0.36,3.95) rectangle ++(0.36,0.36);
\node at (6.45,4.13) {$\cdots$};
\draw[black!60,fill=white] (6.75,3.95) rectangle ++(0.36,0.36);
\draw[arr] (enc.north) |- (5.2,4.13);
\node[font=\scriptsize,text=red!50!black,align=center] at (9.95,4.72) {added input: 48 slots of 6 blank positions,\\to be filled in with $x$, $y$, size and 3 label tokens};
\foreach \b in {0,1} {
  \foreach \i/\t in {0/?,1/?,2/?,3/?,4/?,5/?}
    \draw[red!70!black,fill=red!12] (7.6+\b*2.6+\i*0.36,3.95) rectangle ++(0.36,0.36) node[pos=0.5,font=\scriptsize] {\t};
}
\node at (9.98,4.13) {$\cdots$};
% decoder with the added adapters inside it
\draw[black!60,fill=white,rounded corners=2pt] (5.1,1.9) rectangle (12.5,3.45);
\node[align=center] at (6.25,2.67) {text decoder\\24 layers};
\draw[red!70!black,fill=red!12,rounded corners=2pt] (7.5,2.05) rectangle (12.35,3.3);
\node[align=center,text=red!50!black,font=\scriptsize] at (9.92,2.67) {LoRA adapters on qkv, proj, fc1, fc2 (added)\\in every layer, applied at slot positions only};
\draw[arr] (6.0,3.95) -- (6.0,3.45);
\draw[arr] (9.98,3.95) -- (9.98,3.3);
% readout heads
\node[added,font=\scriptsize] (obj) at (6.3,1.0) {objectness\\(added)};
\node[base,font=\scriptsize] (coord) at (8.55,1.0) {Moondream's coordinate\\and size decoders};
\node[base,font=\scriptsize] (vocab) at (11.0,1.0) {Moondream's\\vocabulary head};
\draw[arr] (7.0,1.9) |- (obj.east);
\draw[arr] (8.55,1.9) -- (coord.north);
\draw[arr] (11.0,1.9) -- (vocab.north);
\draw[black!70,thick] (obj.south) -- (6.3,0.35) -- (11.0,0.35);
\draw[black!70,thick] (coord.south) -- (8.55,0.35);
\draw[black!70,thick] (vocab.south) -- (11.0,0.35) -- (13.1,0.35);
% ---------------- output
\node[font=\bfseries] at (15.2,5.85) {Output: every slot at once};
\node[outbox,text width=4.1cm] (out) at (15.2,3.05) {the same 48 slots, filled in:\\``nothing'', or a box ($x$, $y$, size\\in 1,024 bins), a label and a score\\[3pt]
\textcolor{orange!80!black}{$\blacksquare$} person, 0.82\\
\textcolor{orange!80!black}{$\blacksquare$} cup, 0.41\\
\phantom{$\blacksquare$} nothing\\
\phantom{$\blacksquare$} $\cdots$ (48 slots)};
\draw[arr] (13.1,0.35) -| (out.south -| 13.6,0);
% ---------------- training band
\draw[green!40!black,dashed,thick,rounded corners=4pt] (0.0,-2.55) rectangle (17.5,-0.25);
\node[anchor=north west,font=\bfseries,text=green!30!black] at (0.1,-0.3) {Training: only the added (red) parts receive gradients};
\node[train,text width=4.3cm] (data) at (2.45,-1.6) {labelled objects: COCO train\\(46k images) + 18k images\\labelled by Gemini 4 Carbon};
\node[train,text width=4.4cm] (hung) at (7.6,-1.6) {Hungarian matching: each\\labelled object takes one slot;\\the other slots learn ``nothing''};
\node[train,text width=5.0cm] (loss) at (13.9,-1.6) {losses: box bins (cross-entropy),\\L1 + GIoU, label tokens, objectness\\(an extra denoising pass with some\\true positions revealed adds nothing)};
\draw[arr] (data) -- (hung);
\draw[arr] (hung) -- (loss);
\draw[arr,dashed] (out.south -| 16.4,0) -- (loss.north -| 16.4,0) node[pos=0.5,right,font=\scriptsize] {predictions};
\draw[-{Latex[length=1.8mm]},red!70!black,dashed,thick] (loss.north -| 12.95,0) -- (12.95,2.67) -- (12.35,2.67);
\node[font=\scriptsize,text=red!60!black,anchor=west] at (13.0,-0.05) {gradients};
% legend
\node[base,font=\scriptsize] at (2.6,-3.0) {frozen base model};
\node[added,font=\scriptsize] at (5.3,-3.0) {added and trained};
\node[outbox,font=\scriptsize] at (7.5,-3.0) {output};
\node[train,font=\scriptsize] at (9.35,-3.0) {training only};
\end{tikzpicture}}
\caption{The base model, what we add, what comes out and how it is trained. Moondream~2
(grey) is unchanged: its vision encoder, text decoder and its own coordinate, size and
vocabulary heads are frozen. We add (red) 48 slots of six blank positions after the
image tokens, whose outputs the model fills in (each blank's input is a learned slot query, role and mask embedding),
low-rank adapters that change the decoder only at those positions, and an objectness
layer: about 50M trained parameters at rank 64. One forward pass reads every slot as
``nothing'' or a box, label and score. In training, labelled objects are matched to
slots and the losses update only the added parts.}
\label{fig:architecture}
\end{figure*}

\section{Method}
\textbf{Slots} (Fig.~\ref{fig:architecture}). After Moondream~2's 730-position image
prefix we append 48 slots of six positions. Each position starts as a blank: a learned
mask embedding plus its slot's query and its role (x, y, size, or label token). Positions
see the image and each other, and one forward pass fills them all: positions 0--2 are
read through the region head's 1{,}024-bin coordinate and size decoders (box centre,
width and height), positions 3--5 through the vocabulary head (a three-token label), and
position~0 also gives an objectness logit. This is parallel masked-token decoding as in
MaskGIT~\citep{chang2022maskgit} and masked diffusion
models~\citep{austin2021d3pm,sahoo2024mdlm,nie2025llada}, with one step.

\textbf{Training.} The backbone stays frozen. LoRA adapters~\citep{hu2022lora} of rank
16--128 act only at slot positions, so the image prefix and the model's existing skills
are untouched. Each image's truth set is matched to slots by the Hungarian
algorithm~\citep{kuhn1955hungarian,carion2020detr} on box distance, overlap, objectness
and label likelihood. Matched slots receive cross-entropy on the four bins, L1 and
GIoU~\citep{rezatofighi2019giou} on the expected-bin box, and label cross-entropy;
unmatched slots are pushed to ``nothing'' with weight 0.1. The reported heads also train
a second, denoising pass per step that reveals a random share of each matched slot's
true positions; training without it gives the same results in three quarters of the time
(Table~\ref{tab:ablate}). Runs use one H100, batch~16, 20{,}000 steps (2--4\,h).

\textbf{Selection.} A slot's score is objectness $\times$ label confidence $\times$
(box confidence)$^{1/4}$. Same-label boxes overlapping by IoU 0.5 are fused into their
score-weighted mean~\citep{solovyev2021wbf}. Optionally the head also reads a $2\times2$
grid of tiles overlapping by a quarter, in one batch with the full image, and keeps tile
boxes smaller than four grid cells that no tile edge cuts~\citep{akyon2022sahi}. Every
choice and threshold is tuned on 500 further COCO val images disjoint from the test set.

\textbf{Data.} 46k COCO train2017 images~\citep{lin2014coco} (the 80 categories), plus
images from Open Images~\citep{kuznetsova2020openimages} and
Objects365~\citep{shao2019objects365} labelled by a VLM teacher~\citep{hinton2015distilling}.
A pilot scored twelve VLMs, called through OpenRouter~\citep{openrouter}, against COCO's
labels on 200 val images. Gemini~4 Carbon~\citep{gemini2023} labelled best (F1 0.756,
as complete as Moondream~3.1 and more precise) at \$0.005 per image, so it labelled 21k
images for \$120, of which 17.7k passed the label filter.

\section{Results}
We evaluate on a seeded sample of 500 COCO val2017 images (1{,}446 phrases: every
category present is prompted) and of 1{,}000 Pascal VOC 2007 test images (1{,}548
phrases, the 20 categories under their COCO names), with micro precision, recall and F1
under one-to-one matching. Moondream is prompted once per present category and gives no
confidences, so it has one operating point; the slot head receives no prompt.

\begin{table*}[t]
\centering\small\setlength{\tabcolsep}{4pt}
\caption{COCO val2017 (500 images) and VOC 2007 test (1{,}000 images). Seconds per
COCO image on one L40S (2.9 categories per image); last column: per video frame with all
80 categories, detection every third frame. Slot heads use thresholds tuned on COCO.}
\label{tab:main}
\begin{tabular}{@{}lcccccrr@{}}
\toprule
Model & COCO P & COCO R & COCO F1$_{.5}$ & COCO F1$_{.75}$ & VOC F1$_{.5}$ & s/img & s/frame \\
\midrule
Moondream 2 \texttt{detect} (1.9B) & .500 & \textbf{.709} & .587 & .445 & .851 & .220 & -- \\
Moondream 3.1 (9.3B MoE) & .682 & \textbf{.719} & .700 & .560 & .857 & .208 & 2.79 \\
\midrule
Slots r128, COCO + teacher, fast & \textbf{.811} & .626 & .707 & .566 & .858 & \textbf{.091} & \textbf{.047} \\
Slots r128, COCO + teacher, tiled & .743 & .695 & \textbf{.718} & \textbf{.586} & .844 & .229 & -- \\
Slots r64, COCO only, fast & .809 & .608 & .694 & .559 & \textbf{.862} & .087 & -- \\
\bottomrule
\end{tabular}
\end{table*}

\textbf{Similar quality.} The fast head is level with Moondream~3.1 on COCO and VOC
(Table~\ref{tab:main}) and gains 12 points of COCO F1 over its own backbone's
\texttt{detect}; it trades recall for precision. Paired bootstrap intervals put
differences of about 1.5 points within noise. Moondream~3.1's COCO numbers do not
reproduce its published ones under an unspecified protocol.

\textbf{Small objects.} All models read the same $27\times27$ feature grid. On objects
of four cells or more our recall equals Moondream's (0.87); below that it falls behind,
to 0.22 against 0.36 for objects under one cell, although Moondream finds them in the
same features. Tiling closes that gap (0.37) and lifts COCO F1 to 0.718, at 2.5 times the
time; on VOC, whose objects are larger, it costs precision.

\textbf{Much faster.} The slot pass costs 0.016\,s; the vision encoder dominates the
0.09\,s per image. Cost no longer depends on categories or objects: with all 80 COCO
categories requested, the head needs 0.047\,s per video frame against 2.79\,s for
Moondream~3.1, about 59$\times$ less.

\begin{figure}[t]
\centering
\includegraphics[width=0.9\columnwidth]{precision-recall.png}
\caption{Precision against recall at IoU 0.5 on the 500 COCO test images, before the
selection changes. Dots: each head at the threshold tuned on the disjoint split.
Diamonds: Moondream~2 and 3.1, which give no confidence.}
\label{fig:pr}
\end{figure}

\textbf{What worked and what did not} (Table~\ref{tab:ablate}). Reading boxes through
the model's own coordinate bins instead of regressing four numbers was the largest gain.
Data, steps, keeping small objects, adapter rank 16$\rightarrow$64, adding
teacher-labelled images to COCO, the combined score and tiling each helped. The rest did
not move COCO beyond noise: rank 128, more decoding passes, training without the
denoising pass, category prompts in front of the slots, size-weighted losses and
human Objects365 labels in place of the teacher's (which did lift VOC precision). Slot
groups by object size, ``nothing'' targets only near labelled objects, objectness
trained toward IoU, teacher labels \emph{instead of} COCO's, and mean-teacher
self-training~\citep{tarvainen2017meanteacher,liu2021unbiasedteacher,xie2020noisystudent}
hurt. Custom fused CUDA kernels for Moondream~2's sequential decoder cut its per-box cost
1.85$\times$, but are not bit-exact with eager PyTorch.

\begin{table}[t]
\centering\small\setlength{\tabcolsep}{3pt}
\caption{One change at a time, COCO val F1 at IoU 0.5 (before $\rightarrow$ after); later rows use the selection of Section 2.}
\label{tab:ablate}
\begin{tabular}{@{}clr@{}}
\toprule
& Change & F1$_{.5}$ \\
\midrule
\ok & Coordinate bins instead of regressed boxes (11k img.) & .429$\rightarrow$.563 \\
\ok & 4$\times$ data, 3$\times$ steps & .563$\rightarrow$.622 \\
\ok & Keep small objects (small recall .10$\rightarrow$.26) & .622$\rightarrow$.636 \\
\ok & LoRA rank 16$\rightarrow$64 & .636$\rightarrow$.657 \\
\ok & + 18k teacher-labelled images & .657$\rightarrow$.673 \\
\ok & Combined score, duplicate removal & .674$\rightarrow$.694 \\
\ok & Fusion at IoU 0.5, box confidence$^{1/4}$ & .700$\rightarrow$.707 \\
\ok & $2\times2$ tiles (recall $<$1 cell .21$\rightarrow$.37) & .707$\rightarrow$.718 \\
\eq & LoRA rank 64$\rightarrow$128 & .694$\rightarrow$.700 \\
\eq & 2--3 decoding passes instead of 1 & unchanged \\
\eq & No denoising pass in training & .705$\rightarrow$.705 \\
\eq & Category prompts before the slots & .705$\rightarrow$.702 \\
\eq & Size-weighted loss (recall $<$1 cell +.03) & .705$\rightarrow$.700 \\
\eq & Objects365 human labels for teacher's (VOC +.013) & .705$\rightarrow$.698 \\
\no & Two slot groups by object size (23k img.) & .590$\rightarrow$.549 \\
\no & Objectness target = matched IoU (small recall $-$.07) & .636$\rightarrow$.631 \\
\no & ``Nothing'' only near labelled objects & .636$\rightarrow$.568 \\
\no & Teacher labels only, no COCO (17.7k vs.\ 46k img.) & .694$\rightarrow$.681 \\
\no & EMA self-training, ignore confident slots & .694$\rightarrow$.685 \\
\bottomrule
\end{tabular}
\end{table}

\section{Limitations and conclusion}
The slot head is trained on COCO train images and knows COCO's 80 categories, whereas
Moondream answers free-form phrases; Moondream probably saw COCO in training too, so
VOC, where both are out of domain, is the fairer comparison. We have not evaluated on
LVIS~\citep{gupta2019lvis} or referring expressions~\citep{yu2016refcoco}. Within that
scope, a frozen 1.9B VLM with a 50--100M-parameter slot adapter detects all categories
in one pass at the quality of a 9B model in less than half its time, and in a small
fraction of it when many categories are asked for. Open problems are recall on small
objects without tiling, open-vocabulary labels, where multi-step decoding may matter,
and a cheaper vision encoder.

\bibliographystyle{plainnat}\bibliography{references}
\end{document}
